\scanpage{43}
% Pencil archive foliation: 42.
\deleted{
si le \unsure{fais.} motif de coefficient $\mathbb{Z}(0)$
est remplacé par un motif constant
tordu de \inserted{pur} type dimensionnel \deleted{de} $\delta$
[Attention au décalage par $\delta$ !]. Cet
énoncé \deleted{est un} \uline{\unsure{précise|précède}} les interprétations
galoisiennes et la conjecture de Tate
des \unsure{problèmes}, et est sans doute moins
profond.
}
\margin{\deleted{
$X$ de dim $n$, $U$ ouvert, $X-U$ de dim $<n$.
Alors si \unsure{$i\leq n$}, $H^i(X)\deleted{\unsure{\simeq}}\unsure{\longrightarrow|-}H^i(U)$
est un isom \uuline{mod}
$D_{\unsure{i-1}}(\mathcal{M}^+(K))$, donc $H^i(X)$ et
$H^i(U)$ ont \unsure{$=$|$\overline{\mathrm m}$} composant dans
$D_{\unsure{i|n}}(\mathcal{M}^+(K))/D_{\unsure{i-1|n-1}}(\mathcal{M}^+(K))$.
De plus (via résolution) ce dernier est dans
l'image de $\mathcal{M}_i^+(K)$, d'où un inv. birationnel
$\in\mathbb{Z}^{\unsure{\sigma_i}}\ldots$
}}

\unsure{12)} \deleted{\unsure{B.}}
\uline{L'hom fondamental}
\[
L(K)\unsure{-|\longrightarrow}\mathrm{M}^+(K)
\]
\uline{et invariants birationnels fondamentaux}.

$K$ est un corps. On désigne par
$L(K)$ le gpe \unsure{$K$} engendré par les
préschémas de type fini sur $K$, avec
comme relations celles définies par
découpage en morceaux. C'est un
anneau par $\otimes$. On définit
un hom d'anneaux
\[
\sigma:L(K)\longrightarrow\mathrm{M}^+(K)
\]
